% ======================================================= REPRODUCIBILITY ==
\section{Reproducibility}\label{sec:repro}

\subsection{Environment and hardware}
All production runs used CPU only on one AWS EC2 \code{r7i.2xlarge} instance (8 vCPU, 64 GiB RAM, Ubuntu with glibc
2.39) and Python 3.12.3 with the pinned packages in \cref{tab:env}. No GPU is needed. The clean regeneration of TF3
took 1 h 33 min with six threads, including rebuilding every cache from the raw release. Regenerating ES4 is estimated
at about 3 hours or less and both files at about 4.5 hours or less; neither total has been measured as one run. Peak
memory was not recorded, so 64 GiB is what was used and not a proven minimum. The release is 7.76 GiB as a zip file, its
dense caches take 4.1 GiB, and we recommend at least 40 GiB of free disk.

\paragraph{Kaggle Notebook} The method cannot be run inside a Kaggle Notebook as packaged and was never run in one. It
was developed and operated as a chain of command-line stages under an exact pinned environment on a larger machine than
a standard notebook provides. We therefore supply a repository and a written account of how it was run,
\code{REPRODUCIBILITY.md}, which lists the versions, preprocessing steps, seeds, commands and verification status.

\begin{table}[h]
\centering\small
\caption{Pinned inference environment (\code{requirements.txt}) and the licence of each dependency.}\label{tab:env}
\begin{tabular}{@{}l l l@{}}
\toprule
Package & Version & Licence \\
\midrule
NumPy \cite{harris2020} & 2.5.3 & BSD-3-Clause \\
pandas & 3.0.6 & BSD-3-Clause \\
SciPy \cite{virtanen2020} & 1.18.1 & BSD-3-Clause \\
PyArrow & 25.0.1 & Apache-2.0 \\
LightGBM \cite{ke2017} & 4.7.0 & MIT \\
scikit-learn & 1.9.1 & BSD-3-Clause \\
\bottomrule
\end{tabular}
\end{table}

\subsection{Commands}
The repository has one entry point. The competition release is not redistributed; \code{-{}-release-root} is the
directory that contains \code{submission\_key.csv}, \code{corridors/} and \code{task2/}.

\smallskip
\noindent\begin{minipage}{\linewidth}
\small
\begin{verbatim}
# 1. verify the final archives and the TF3 code bundle (standard library only)
python reproduce.py

# 2. regenerate TF3 from the raw release and the saved weights
python -m pip install -r requirements.txt
python reproduce.py --submission tf3 --release-root /path/to/release \
                    --out /path/to/new_tf3 --threads 6

# 3. regenerate ES4 from the raw release, including EQ1 retraining
python reproduce.py --submission es4 --release-root /path/to/release \
                    --out /path/to/new_es4 --threads 6

# 4. regenerate both selected files, sharing one set of caches
python reproduce.py --submission all --release-root /path/to/release \
                    --out /path/to/new_all --threads 6
\end{verbatim}
\end{minipage}

\subsection{What has been verified}
\begin{itemize}
  \item \textbf{TF3, end to end.} The bundled driver rebuilds every cache and prediction component from the release
        and the saved TRAIN weights, then compares all 6,980,503 rows with the submitted file. On the clean run every queue
        label matched exactly and the largest continuous difference was $1.82\times10^{-12}$. This verifies saved-weight
        inference; the refit procedure, seeds and inventories are recorded, but an independent retraining was not run.
  \item \textbf{ES4.} The final archive is in the repository and its checksum is verified by \code{reproduce.py}.
        Row-level comparison confirms that ES4 is TF3 plus 247 ongoing-queue labels and 5,825,667 recentered speeds.
        \code{code/regenerate\_es4.py} runs the original stage scripts from the release through ES1, ES2, ES3 and ES4 and
        compares every row with the frozen file. It has been exercised on synthetic data but \textbf{has not yet been run
        on the real release}. The EQ1 classifier weights existed only on the instance used during the competition, which
        no longer exists, so the command retrains EQ1 from the release; whether that model is identical to the original
        is not yet known. Only the 247 EQ1-dependent labels depend on this. TF3 scored 0.00196 (0.22\,\%) below ES4 on the
        validation month.
\end{itemize}

\subsection{Artefacts and checksums}
\begin{table}[H]
\centering\small
\caption{SHA256 of the main artefacts. The first three are verified by \code{python reproduce.py}.}\label{tab:sha}
\begin{tabular}{@{}l l@{}}
\toprule
Artefact & SHA256 \\
\midrule
\code{submissions/submission\_TF3.zip} & \sha{2ae94d422399c7248297487fc8f45d737de0e024fbad32b6a360c290bbabca7a} \\
\code{submissions/submission\_ES4.zip} & \sha{ce03386b42e33bc6c726b09c0b64b791f294ea3abeb88d7e310552bb080c04bb} \\
\code{vendor/reproduction\_TF3.zip} & \sha{dce93b3f2a38d27014497ec1a563cc1e2d032da11c499234f226c72093678cd7} \\
TS3 frozen rule & \sha{61499049f1cced66927eff4c9e9c99bb611093c33c3a31de98ebd436940883d5} \\
TS3 final model manifest & \sha{8ed505b54215cd88706e95dedc5b4fbee3f435e228560f8d71d018b546cf6acb} \\
TS3 registered protocol & \sha{fcc9f299f32bc0df8dc68749df21ceb1d2bc2609fc3d2a21b15a85913c84aec0} \\
Q46 ongoing classifier (TF3) & \sha{2816e1d3f3c233416327591c6659cd9e80a25f0054431f1b1d4ee23735b2ab46} \\
EQ1 ongoing classifier, 800 trees (ES4) & \sha{cb626813e621064594f6956344df81d7d8879e65ebc23ae2f164a2f8f6f9ac9f} \\
\bottomrule
\end{tabular}
\end{table}

\subsection{Seeds and data membership}
TS3 members draw their training rows with \code{random\_state} equal to the member seed (20261104 or 20261204) plus
the panel index; the RS1 models use LightGBM seed 20261001 and samples drawn with 20261001 plus the panel index; Q46's
training windows use 20261046 plus the panel index; EQ1's additional windows use 20261240 plus the panel index; both
classifiers use LightGBM seed 20261002. At inference, the only seed is the S10 randomized start, $20261003 + 100\times$
panel index $+$ month. The fit manifests in the TF3 bundle record the fitting months, row counts and input
hashes of every saved model, and two independent audits recomputed every stored weekday profile (31,783,763 fitting
rows) exactly.

\subsection{External data and licence declaration}
No external dataset, pretrained model, paid or rate-limited API or proprietary dependency was used. The only data are
the competition release, used under the competition rules and not redistributed here. All software dependencies are
open source under the licences in \cref{tab:env}. Figures were produced with Matplotlib and the report with \LaTeX.
Our own code is released under the MIT licence.
